\honest{The True Prisoner's Dilemma}{Eliezer Yudkowsky}{2008}

It occurred to me one day, I say, that the standard way of picturing the Prisoner's Dilemma is fake.

The game is a table of payoffs. You get 5 for defecting while the other cooperates, 3 if you both cooperate, 2 if you both defect, and 0 for cooperating while the other defects. Whatever the other does, defecting pays more, so you both defect, though you would both rather have cooperated.

The usual story has two captured criminals deciding whether to testify against each other, or two strangers playing once for dollars. And, oh yes, you are supposed to ``pretend that you're entirely selfish.'' That, in my view, is what makes it fake. \nb{The stipulation is standard: William Poundstone's version, as Wikipedia quotes it, says ``Each prisoner is concerned only with his own welfare.'' Game theory already counts everything a player cares about in the payoffs. Robin Hanson said so in the comments, and Yudkowsky replied that the complaint was about how the game is illustrated.}

You cannot avoid hindsight bias by telling a jury to ignore the outcome, and without a complicated effort a normal human being cannot pretend to be truly selfish. \nb{The first is a finding: instructing mock jurors made no difference. The second rests on the comparison alone.} We are born with fairness, sympathy and even altruism, because our ancestors played the iterated dilemma. \nb{This is Robert Trivers's theory of reciprocal altruism (1971), which lists sympathy among the traits it explains.} So we do not entirely prefer exploiting a cooperator to cooperating. In the laboratory, I say, we are not entirely unsympathetic to the stranger, and our instinct is to ask how to get mutual cooperation. \nb{No experiment is cited. In one well-known public-goods experiment, half the subjects gave more when others gave more, and about a third free-rode completely.} For such a player, mutual cooperation beats exploiting the other, and the question is whether the other sees it the same way. \nb{Game theorists call a game with that order a stag hunt, or an assurance game. In 1978 Kelley and Thibaut had described how a person's concerns turn the material outcomes into the matrix they act on.}

So I build a true dilemma. Four billion people are dying of a disease that only substance S can cure, and S must be shared with a paperclip maximizer from another dimension that we will never meet again. Mutual cooperation saves 2 billion lives and makes 2 paperclips; if only we defect, 3 billion lives and no paperclips; if only it defects, no lives and 3 paperclips; mutual defection, 1 billion and 1. I chose the numbers to make you indignant. Here we really do prefer defecting against a cooperator, and both sides still prefer mutual cooperation to mutual defection.

I end by asking what you would do. \nb{The post gives no answer. In the comments Yudkowsky wrote ``I didn't say I would defect,'' and the next day posted ``The Truly Iterated Prisoner's Dilemma'', which is not in this book.}
